\documentclass[10pt,a4paper]{amsart}
\usepackage[margin=19mm]{geometry}
\usepackage{amsmath,amssymb,mathtools}
\usepackage[scr=rsfs,cal=euler]{mathalfa}
\usepackage{xfrac}
\usepackage[unicode,hidelinks,bookmarks=false]{hyperref}
\newcommand{\ie}{i.e.\ }
\newcommand{\cf}{cf.\ }
\newcommand{\resp}{respectively\ }
\newcommand{\Subsection}{\subsection}
\providecommand{\nobreakdash}{\nobreak\hbox{-}}
\setlength{\emergencystretch}{2em}
\multlinegap0pt
\usepackage{mathtools}
\usepackage{amssymb}
\newtheorem{prop}{Proposition}
\newcommand{\FY}{\operatorname{FY}}
\newcommand{\Hilb}{{\mathsf{Hilb}}}
\renewcommand{\L}{{\mathcal{L}}}
\newcommand{\rk}{{\mathrm{rk}}}
\title{Equivariant gamma-positivity of matroid Chow rings}
\author{Hsin-Chieh Liao}
\date{Source: arXiv:2408.00745v3 (CC BY 4.0)}
\begin{document}
\maketitle

\noindent\emph{Source and licence.} Equivariant gamma-positivity of matroid Chow rings. Version arXiv:2408.00745v3 (CC BY 4.0), distributed by arXiv under the Creative Commons Attribution 4.0 International licence, http://creativecommons.org/licenses/by/4.0/, as recorded in the arXiv licence metadata for this exact version and shown on the abstract page https://arxiv.org/abs/2408.00745v3. Full licence notice: \texttt{licenses/arxiv-2408.00745-CC-BY-4.0.txt}.\par\smallskip

\noindent\textit{Editorial scope.} Counted unit: the article's prop prop:HilbChow (source0.tex lines 321--347). The article's complete original proof of it is delivered with it (source0.tex lines 348--403, one \texttt{proof} environment of 56 lines). No mathematics is written, repaired or re-derived: the statement and the proof are the article's own lines, in the article's own notation, and no retained line is substituted or re-flowed.  No further source-local context is needed: every cross-reference of every retained line resolves inside the two delivered spans.  Nothing was omitted from any retained span, so the package carries no omission record.  No retained line contains a citation, so the wrapper carries no bibliography and no bibliography entry.  No retained line contains a figure environment, an image-inclusion command or drawing code, so no image is delivered and the package declares no asset dependency.  Editorial formatting only, in addition: an AMS \texttt{amsart} preamble that restates the article's own declarations and macros used by the retained lines, namely the environments \texttt{prop} on separate counters, together with the standard packages those lines need.  The retained environments print in this document as Proposition 1 in order of appearance, whereas the article prints the same items with the numbering of its own full document.  The complete native source of the article as distributed in the arXiv e-print for this exact version is bundled unchanged as \texttt{sources/163228/source0.tex}.
\begin{prop} \label{prop:HilbChow}
Let $M$ be a loopless matroid of rank $r$.
\begin{itemize} 
    \item[(i)] The $G$-equivariant Hilbert series of $A(M)_{\mathbb{C}}$ is given by
    \[
        \Hilb_G(A(M)_{\mathbb{C}},t)=\sum_{S\subseteq [r-1]}\phi_{S,r}(t)[\alpha_{\L(M)}(S)],
    \]
    where for a subset $S=\{s_1<s_2<\cdots<s_{\ell}\}\subseteq [r-1]$,
    \[
        \phi_{S,r}(t)\coloneqq t^{|S|}[s_1-1]_t[s_2-s_1-1]_t\cdots[s_{\ell}-s_{\ell-1}-1]_t[r-s_{\ell}]_t.
    \]
    In particular, $\phi_{\emptyset,r}(t)=[r]_t$.
    \item[(ii)] The $G$-equivariant Hilbert series of $\widetilde{A}(M)_{\mathbb{C}}$ is given by
    \[
    \Hilb_G(\widetilde{A}(M)_{\mathbb{C}},t)=\sum_{S\subseteq[r-1]}\psi_{S,r}(t)[\alpha_{\L(M)}(S)],
    \]
    where for a subset $S=\{s_1<s_2<\cdots<s_{\ell}\}\subseteq [r-1]$,
    \[
    \psi_{S,r}(t)=\begin{cases}
        t^{|S|}[s_1]_t[s_2-s_1-1]_t\ldots[s_{\ell}-s_{\ell-1}-1]_t[r-s_{\ell}]_t & \text{ if }S\neq \emptyset\\
        [r+1]_t & \text{ if } S=\emptyset
    \end{cases}.
    \]
\end{itemize}


\end{prop}

\begin{proof}
For (i), recall that the Feichtenr-Yuzvinsky basis $\FY(M)$ consists of monomials $x_{F_1}^{a_1}x_{F_2}^{a_2}\ldots x_{F_{\ell}}^{a_{\ell}}$ indexed by chains (including the empty chain)
\[
    \emptyset\neq F_1\subsetneq F_2\subsetneq\cdots\subsetneq F_{\ell}\subseteq{E}
\]
such that the exponent $a_j$ satisfies $1\le a_j\le \rk(F_j)-\rk(F_{j-1})-1$ for $j=1,2,\ldots,\ell$, where we set $F_0=\emptyset$.
For any subset $S=\{s_1<s_2<\cdots<s_{\ell}\}\subseteq [r-1]$, the permutation $\mathbb{C}G$-module generated by chains in $\L(M)$ with rank set $S$ is
\begin{align*}
    \alpha_{\L(M)}(S)
    &=\mathbb{C}G\{F_1\subsetneq \cdots\subsetneq F_{\ell}: \rk(F_i)=s_i~\forall i\}.\\
    &\cong \mathbb{C}G\{F_1\subsetneq \cdots\subsetneq F_{\ell}\subsetneq E: \rk(F_i)=s_i~\forall i\}.
\end{align*}
Hence, as a graded $\mathbb{C}G$-module, $A(M)_{\mathbb{C}}$ decomposes as a direct sum of $\alpha_{\L(M)}(S)$ over all $S\subseteq [r-1]$.

Assume that $
    \Hilb_G(A(M)_{\mathbb{C}},t)=\sum_{S\subseteq [r-1]}\phi_{S,r}(t)[\alpha_{\L(M)}(S)]$
for some polynomials $\phi_{S,r}(t)$. To determine $\phi_{S,t}(t)$, we consider a map $f:\FY(M)\longrightarrow 2^{[r-1]}$ defined by 
\[
    f(x_{F_1}^{a_1}x_{F_2}^{a_2}\ldots x_{F_{\ell}}^{a_{\ell}}x_E^i)=\{\rk(F_1),\rk(F_2),\ldots,\rk(F_{\ell})\}
\]
where all flats $F_i$ are distinct and are not $E$. Since for any $S=\{s_1<\cdots<s_{\ell}\}\subseteq [r-1]$, the inverse image of $f$ on $S$ consists of monomials
\[
    f^{-1}(S)=\left\{x_{F_1}^{a_1}x_{F_2}^{a_2}\ldots x_{F_{\ell}}^{a_{\ell}}x_E^i : \substack{\emptyset=F_0\subsetneq F_1\subsetneq\cdots\subsetneq F_{\ell}\subsetneq E \text{ with }\rk(F_j)=s_j\\
    1\le a_j\le s_j-s_{j-1}-1  \text{ for }j=1,2\ldots,\ell\\
    0\le i\le r-s_{\ell}-1
    }\right\},
\]
Therefore,
\begin{align*}
    \phi_{S,r}(t)
    &=(t+t^2+\cdots+t^{s_1-1})\ldots(t+t^2+\cdots+t^{s_{\ell}-s_{\ell-1}-1})(1+t+\cdots+t^{r-s_{\ell}-1})\\
    &=t^{|S|}[s_1-1]_t[s_2-s_1-1]_t\ldots[s_{\ell}-s_{\ell-1}-1]_t[r-s_{\ell}]_t.
\end{align*}
In particular, when $S$ is an empty set, the inverse image $f^{-1}(\emptyset)=\{x_E^i~:~0\le i\le r-1\}$ and $\phi_{\emptyset,r}(t)=[r]_t$.

For (ii), the argument is analogous. The basis $\widetilde{\FY}(M)$ consists of monomials indexed by chains as above, with exponents satisfying
$1\le a_1\le \rk(F_1)$ and $1\le a_j\le \rk(F_j)-\rk(F_{j-1})-1$ for $j\ge 2$. 
Thus, we write
\[
\Hilb_G(\widetilde{A}(M)_{\mathbb{C}},t)
=\sum_{S\subseteq [r-1]} \psi_{S,r}(t)\,[\alpha_{\L(M)}(S)].
\]
Define $g:\widetilde{\FY}(M)\longrightarrow 2^{[r-1]}$ by
\[
g(x_{F_1}^{a_1}\cdots x_{F_\ell}^{a_\ell}x_E^i)
=\{\rk(F_1),\ldots,\rk(F_\ell)\}.
\]
Then counting $g^{-1}(S)$ for $S=\{s_1<\ldots<s_{\ell}\}\subseteq [r-1]$ gives 
\[
    \psi_{S,r}(t)=\begin{cases}
        t^{|S|}[s_1]_t[s_2-s_1-1]_t\cdots[s_{\ell}-s_{\ell-1}-1]_t[r-s_{\ell}]_t & \text{ if }S\neq \emptyset\\
        [r+1]_t & \text{ if } S=\emptyset
    \end{cases}.,
\]
as claimed.
\end{proof}

\end{document}
